\section{Paired electric and magnetic radial control}
\label{cd:section}

The distant finite-order control construction can be tracked in both
polarizations simultaneously. Its limiting TE and TM transfers obey a
specific inverse-transpose relation. This supplies controlled amplification
of an electric-wall reflection sector and suppression of its opposite
sector, within any prescribed finite band. It does not bound unselected
partial waves or the Green operator of a perturbed device.

\subsection{The same physical layers control a paired transfer graph}

Retain $X,Y,J,q,\ell_M,b_M,d_M$ from Section~\ref{cc:section}.
For TE use $v_{\rm E}=(u,u')^T$ and for TM use
$v_{\rm M}=(u,n^{-2}u')^T$ in an index-$n$ layer. With
$\nu=l(l+1)$ their exact dimensionless equations are
\begin{equation}\label{cd:generators}
 v_{\rm E}'=(X-n^2Y+\nu x^{-2}Y)v_{\rm E},\qquad
 v_{\rm M}'=(n^2X-Y+\nu n^{-2}x^{-2}Y)v_{\rm M}.
\end{equation}
Both states are continuous across each material interface.

\begin{lemma}[PROVED: paired angular derivatives of the identity loop]
\label{cd:loop}
For the high--low, then low--high loop of
Lemma~\ref{cc:loop-derivative}, its two reference propagators satisfy
$U_{0,{\rm M}}=U_{0,{\rm E}}^{-T}$ at every prefix.
Their integrated angular derivatives are
\begin{equation}\label{cd:loop-derivatives}
 B_{{\rm E},M}=\begin{pmatrix}0&-d_M\\b_M&0\end{pmatrix},
 \qquad
 B_{{\rm M},M}=\begin{pmatrix}0&-b_M\\d_M&0\end{pmatrix}
                 =-B_{{\rm E},M}^{T}.
\end{equation}
For both polarizations the integrated interaction-generator norm is
$b_M+d_M\leq2b_M$. Consequently the one-loop error
\eqref{cc:one-loop-error} and the repeated-loop error
\eqref{cc:global-error} apply simultaneously, with the same constants.
The limiting repeated-loop transfers are
\begin{equation}\label{cd:shear-pair}
 \bigl(e^{t\nu Y},e^{-t\nu X}\bigr),\qquad t\geq0.
\end{equation}
\end{lemma}
\begin{proof}
The two reference generators are negative transposes:
$n^2X-Y=-(X-n^2Y)^T$. Differentiating an inverse transpose proves
the prefix identity, including across continuous state interfaces.
Write the TE prefix as
$U=\left(\begin{smallmatrix}a&b\\c&d\end{smallmatrix}\right)$.
The TM prefix is
$U^{-T}=\left(\begin{smallmatrix}d&-c\\-b&a\end{smallmatrix}\right)$.
Its angular interaction matrix, including the actual factor $n^{-2}$,
is therefore
\begin{equation}\label{cd:tm-interaction}
 n^{-2}U_{0,{\rm M}}^{-1}YU_{0,{\rm M}}
   =n^{-2}\begin{pmatrix}cd&-c^2\\d^2&-cd\end{pmatrix}.
\end{equation}
Its norm is $(c^2+d^2)/n^2$.

For completeness, in the forward half with $0\leq j<M$, the TE
second rows during the high and low quarters, up to a common sign, are
\[
 (-q^{j+1}\sin\theta,q^{-j}\cos\theta),\qquad
 (-q^{j+1}\cos\theta,-q^{-j-1}\sin\theta).
\]
In the reverse half with $M\geq h\geq1$, the corresponding low
and high rows are
\[
 (-q^h\sin\theta,q^{-h}\cos\theta),\qquad
 (-q^h\cos\theta,-q^{1-h}\sin\theta).
\]
Use $ds=d\theta/q$ in a high quarter and $ds=d\theta$ in a
low quarter. Integrating $c^2/n^2$ gives exactly the same
$b_M$ as the integral of $a^2$ in the TE calculation; integrating
$d^2/n^2$ gives $d_M$. These identities follow on each quarter
by interchanging $\sin^2\theta$ and $\cos^2\theta$.
The integral of $cd/n^2$ in a forward pair is
$(1-q^{-2})/2$, and in a reverse pair its negative.
Thus its integral over the loop is zero, proving
\eqref{cd:loop-derivatives} and the norm assertion.

The TM interaction equation differs from its reference equation
by $\nu(r+s)^{-2}$ times \eqref{cd:tm-interaction}.
Its integrated norm and the error from freezing $(r+s)^{-2}$
are consequently exactly bounded as in Lemma~\ref{cc:one-loop}.
Since $B_{{\rm M},M}/b_M=-X+\kappa_qq^{-2M}Y$,
the repeated-loop proof applies with $Y$ replaced by $-X$.
Both square to zero and have norm one, so every constant in
\eqref{cc:global-error} is unchanged.
\end{proof}

\begin{theorem}[PROVED: global finite-order paired TE/TM transfer control]
\label{cd:paired-control}
For distinct positive integers $l_1,\ldots,l_m$ and arbitrary
$G_j\in{\rm SL}(2,\mathbb R)$, there are actual finite radial
dielectrics with $1\leq\varepsilon\leq q^2<\varepsilon_+$ in
$[R,R+d_R]$, $d_R=o(R)$, whose simultaneous transfers converge to
\begin{equation}\label{cd:paired-target}
 (G_j,G_j^{-T})\qquad(1\leq j\leq m)
\end{equation}
as $R\to\infty$. The first component is TE and the second TM.
This is a constructed family of paired targets, not a universal
identity for arbitrary physical radial structures.
\end{theorem}
\begin{proof}
Define asymptotic realizability for the $2m$ transfers together,
using the maximum of their matrix norms. Concatenation and
closedness follow by exactly the threshold-selection proof of
Lemma~\ref{cc:common-and-closure}.
Every common reference transfer $G$ in that lemma is now paired
with $G^{-T}$, since the reference generators are negative
transposes. For a fixed layer word both angular errors are
$O(R^{-2})$, so the common paired tuple is realizable.
Lemma~\ref{cd:loop} supplies the weighted shear pair
\eqref{cd:shear-pair}.

The map $\iota:G\mapsto(G,G^{-T})$ is an injective continuous
group homomorphism: $(GH)^{-T}=G^{-T}H^{-T}$. Its infinitesimal
map is $A\mapsto(A,-A^T)$. Thus common conjugation, the product
formulas, and the integer-frequency signed-time construction in
Theorem~\ref{cc:global-control} all remain within this graph.
In particular the paired shear inverses are supplied by positive
layers through the same periodic rotation argument.
The Vandermonde and open-subgroup proof applies to the first
components and hence, by $\iota$, to their whole graph.
It supplies every tuple \eqref{cd:paired-target} in the closed
paired realizable semigroup. No step replaces TM behavior by an
independently assumed transfer.
\end{proof}

\subsection{Reflection parity and finite-band isolation}

Let $R_z(x,y,z)=(x,y,-z)$. An electric-wall field means
\begin{equation}\label{cd:electric-wall}
 E(R_zx)=-R_zE(x),\qquad H(R_zx)=R_zH(x).
\end{equation}
It has $E_t=0$ and $H_z=0$ on the plane. The opposite,
magnetic-wall field has the reversed two signs and satisfies
$E_z=0$ and $H_t=0$. Each sector has identically zero normal
Poynting density on the plane, for every coherent combination.

\begin{lemma}[PROVED: spherical-wave reflection sectors]
\label{cd:parities}
For azimuthal index $m$, a TE spherical wave is electric-wall
exactly when $l+m$ is even; a TM wave is electric-wall exactly
when $l+m$ is odd. On the subspace invariant under a rotation by
$\pi$ around the $z$ axis, with the vector rotated as well, only
even $m$ occur. Within this $C_2$-even subspace the electric-wall
sector is precisely TE even $l$ together with TM odd $l$.
\end{lemma}
\begin{proof}
The scalar identity is
$Y_{lm}(\pi-\theta,\varphi)=(-1)^{l+m}Y_{lm}(\theta,\varphi)$.
For $X_{lm}=\widehat r\times\nabla_{S^2}Y_{lm}/\sqrt{l(l+1)}$,
reflection reverses the cross product, giving
$X_{lm}(R_zx)=(-1)^{l+m+1}R_zX_{lm}(x)$.
This is the TE electric-field parity. Taking a curl reverses
that reflection sign, so the TM electric-field parity is
$(-1)^{l+m}$. The Maxwell curl equation gives the opposite
magnetic-field sign in each case. This proves the wall
classification. Rotation by $\pi$ contributes $e^{im\pi}$ after
the vector rotation, so its invariant sector has even $m$.
\end{proof}

\begin{theorem}[PROVED: controlled finite-band wall amplification and suppression]
\label{cd:wall-isolation}
Fix any finite set of angular orders and positive numbers $A_l$.
In the $C_2$-even input band of those orders, there is a sequence
of finite ordinary radial dielectrics with vacuum-core radius
tending to infinity such that, after one common rephasing, the
physical core multipliers converge to
\begin{equation}\label{cd:wall-gains}
 \begin{array}{c|cc}
 &{\rm TE}&{\rm TM}\\ \hline
 l\ {\rm even}&A_l&A_l^{-1}\\
 l\ {\rm odd}&A_l^{-1}&A_l
 \end{array}.
\end{equation}
Thus the electric-wall sector has gain $A_l$ and its opposite
sector gain $A_l^{-1}$, for each controlled order. For a uniform
$A_l=A>0$, the finite-band core operator converges to
$A P_{\rm e}+A^{-1}P_{\rm m}$ after that rephasing.
Convergence of all coherent fields and their derivatives is
uniform on each fixed compact core subset. The wall boundary
identities themselves hold exactly at every stage, by symmetry.
\end{theorem}
\begin{proof}
Take the paired targets
$G_l=\operatorname{diag}(A_l,A_l^{-1})$ and restrict starting
radii to $R=2\pi n$. The regular radial state tends to
$v_l^0=(\sin(-l\pi/2),\cos(-l\pi/2))^T$ in either polarization,
because the inner core is vacuum. For even $l$ this is a signed
second coordinate vector; for odd $l$ it is a signed first one.
Thus the limiting endpoint multiplier on this real state is
$A_l^{-1}$ for TE even and TM odd, and $A_l$ for the opposite
two cases.

At the outer endpoint $X_R$, the incoming radial coefficient
obeys the same Wronskian formula in either polarization. The
exterior and core are both vacuum, so their state conventions
agree there. The Hankel calculation in
Corollary~\ref{cc:core-gains} gives
$e^{-iX_R}A_l^{\rm in}\to c_l/2$ when the endpoint state is
$c_l v_l^0$. Its physical core multiplier, relative to the free
wave with the same physical incoming amplitude, is
$(2A_l^{\rm in})^{-1}$. After multiplication by $e^{iX_R}$ it
therefore tends to $c_l^{-1}$, proving \eqref{cd:wall-gains}.
This relative multiplier is valid for TM as well as TE; their
different fixed impedance factors cancel in the free/actual
comparison.

Lemma~\ref{cd:parities} identifies these gains with the two wall
sectors for even $m$. All azimuthal indices within an order have
the same radial transfer. Finite-band convergence gives uniform
compact field and derivative convergence. Radial symmetry
preserves reflection and $C_2$ sectors exactly for every finite
device, establishing the final assertion.
\end{proof}

\begin{remark}[What the suppression does and does not provide]
\label{cd:scope}
For uniform $A$ the limiting multiplier
$D_A=A P_{\rm e}+A^{-1}P_{\rm m}$ preserves the vacuum plane
flux form pointwise:
\begin{equation}\label{cd:flux-cancel}
 D_A^*Q_z^0(x)D_A=Q_z^0(x)\qquad(z=0).
\end{equation}
Indeed both diagonal wall blocks vanish pointwise, and the two
real factors cancel in the remaining blocks. Amplifying the
electric-wall background alone creates no plane passage.

The theorem supplies a controlled finite set of opposite-parity
TM and TE input responses. It supplies no uniform estimate for
unselected angular orders, a total Green operator, a source
inside a new inclusion, or a geometry after a nearby inclusion
breaks the reflection symmetry. Such estimates are required
before a proposed small-transducer construction can be claimed
to satisfy physical positive passage and absolute bypass.
\end{remark}
